\textbf{Problème 4.1 (Dujella, 2026).} Soit $k \geqslant 2$ un entier. Est-il vrai que l'équation de Pell
\[
x^2 - (k^2+1)y^2 = k^2
\]
admet au plus une solution avec $0<y<k-1$ et $x>0$ ?

\medskip

\textit{Remarque.} Matthews, Robertson et White (Glas. Mat. Ser. III \textbf{48} (2013), 265--289) ont vérifié que l'énoncé est vrai pour $k \leqslant 2^{50}$, tandis que Le et Srinivasan (Glas. Mat. Ser. III \textbf{58} (2023), 59--65) l'ont démontré pour tout $k$ de la forme $k=p^mq^n$, où $p,q$ sont des nombres premiers impairs distincts et $m,n$ des entiers positifs.
